% ==============================================================================
% PDF-1 / Section 06 — Invariance Stress Tests (Skeleton)
% No new primitives. Diagnostic-only. Container structure.
% ==============================================================================

\section{Invariance Stress Tests}
\label{sec:invariance-stress-tests}

\subsection{Purpose of invariance testing}
Invariance stress tests are used to verify that diagnostic outcomes remain stable
under admissible transformations of observation, representation, or projection.
They do not test performance, efficiency, or effectiveness.

The objective is to detect violations of invariants already defined in the KEA and
ETS layers.

\subsection{Classes of admissible perturbations}
This section reserves space for discussing admissible perturbations, such as:
\begin{itemize}
  \item reparameterization of observation without semantic change,
  \item partial loss or aggregation of signals within admissible bounds,
  \item alternative but equivalent projection paths.
\end{itemize}

No new perturbation classes are defined here.

\subsection{Stress test outcomes}
Possible outcomes of invariance stress tests include:
\begin{itemize}
  \item confirmed invariance (diagnostic stability),
  \item detected sensitivity (boundary proximity),
  \item violation leading to STOP region identification.
\end{itemize}

All outcomes are descriptive and diagnostic-only.

\subsection{Relation to falsifiability}
Failure of invariance under admissible conditions is treated as a falsification
signal for the diagnostic setup under the current observation context.
This does not falsify the KEA itself, but constrains its applicability.

\subsection{Non-goals}
\begin{itemize}
  \item No robustness optimization.
  \item No tolerance tuning.
  \item No resilience scoring.
  \item No operational recommendations.
\end{itemize}
